% ============================================================
%  Georg Brandes, Samlede Skrifter. Trettende Bind (Supplementbind).
%  Kjøbenhavn: Gyldendalske Boghandels Forlag (F. Hegel & Søn), 1903.
%  Introduction (1903) and "The Controversy over Faith and Knowledge" (1866--67).
%  TRANSLATION (English), printed pp. 1--42.
%  Translated from transcription.tex (Danish) in this directory, whose
%  header records how its text was established.  Method:
%  ../../../../TRANSLATION-PLAYBOOK.md.  The volume is Brandes's 1903
%  revision of pieces first printed 1865--69; the translation follows the
%  1903 text.
%
%  ---- REGISTER ----
%  Readable modern English.  Brandes's long periods are broken where English
%  will not carry them, but nothing is softened: the polemic is sarcastic,
%  and where the Danish is ironic the English is ironic.
%
%  ---- TERMINOLOGY ----
%  Aligned with texts/brandes/dualismen/translation.tex and the Nielsen and
%  Brøchner translations in this repo, since the pieces belong to the same
%  controversy:
%    Tro / Viden          = faith / knowledge   (Tro og Viden = faith and
%                           knowledge)
%    Erkendelse           = cognition (NOT "knowledge", reserved for Viden)
%    Videnskab            = science (Wissenschaft; never "natural science",
%                           for which Brandes writes Naturvidenskab)
%    Fagvidenskaberne     = the special sciences
%    absolut Uensartethed = absolute heterogeneity
%    Aand / aandelig      = spirit / spiritual;  Fornuft = reason;
%    Forstand             = understanding;  Begreb = concept
%    Samvittighed         = conscience;  Selvansvarlighed = self-responsibility
%    Livsanskuelse        = view of life;  Verdensanskuelse = world view
%    Dannelse             = cultivation;  Folkeaand = national spirit
%    Aabenbaring          = revelation;  Under = miracle
%  Danish work-titles stay in Danish, as Brandes cites them (Grundideernes
%  Logik, Om den gode Vilje, Fædrelandet, Dagbladet, ...).
%
%  ---- CONVENTIONS ----
%  - Quotation marks: the volume's single guillemets ‹...› -> ``...''.
%  - Emphasis: italic is the volume's only emphasis device; every \emph{}
%    span of the transcription is carried over as \emph{}, including the
%    Latin and other foreign phrases, which stay untranslated.
%  - Footnotes \fnmark{*)}{...} at the same anchor.
%  - The ɔ: id-est mark -> "i.e."
%  - Page markers \opage{N} and "% --- p. N ---" at the same joints as in
%    transcription.tex, so the two can be read side by side.
%  - Heads: a book under review keeps its Danish title, with an English
%    gloss in brackets; Brandes's own titles are translated.  The table of
%    contents gives the English of the volume's Indhold titles.
%
%  ---- HOW THE TRANSLATION WAS CHECKED (2026-09-22) ----
%  (1) Mechanical audit (tr_audit.py): per printed page, the same \opage
%  joints, \emph spans, footnotes, heads, rules and paragraph breaks as
%  transcription.tex, and an English/Danish word ratio within 0.95--1.60:
%  42 pages, 0 flagged.  (2) Independent review: two readers who had not
%  translated these pages read every sentence of the Danish against the
%  English for omissions, additions, negations, modals, referents and
%  terminology.  No omission found; 7 corrections applied (plus the spacing
%  of eight em-dashes on pp. 13--19 made uniform), namely
%  `vil kunne finde' (p. 8, was "is bound to find"), `bør' (p. 16, was
%  "must"), a doubled `only' (p. 25), `aandelig' as "intellectual" where the
%  sense is debate (pp. 3, 34), an added `merely' (p. 38), and Brandes's
%  self-quotation on p. 35 brought into line with its source on p. 32.
%
%  ---- DEFECTS: DANISH AS PRINTED, ENGLISH CORRECT ----
%  The transcription is diplomatic and carries every printer's error; the
%  translation renders the SENSE and logs the divergence in a % comment at
%  the site.  The English quotation marks balance throughout.
% ============================================================

\documentclass[12pt,a4paper,openany]{book}
\usepackage[utf8]{inputenc}
\usepackage[T1]{fontenc}
\usepackage{libertinus}
\usepackage{libertinust1math}
\usepackage{textalpha}
\usepackage{graphicx}
\usepackage[english]{babel}
\usepackage[protrusion=true,expansion=true]{microtype}
\usepackage{setspace}
\usepackage[margin=1.2in]{geometry}
\usepackage{enumitem}
\usepackage{fancyhdr}
\setlength{\headheight}{28pt}

% Footnotes as the 1903 printing sets them: *), set at each site with
% \fnmark{*)}{...}, exactly as in transcription.tex.
\renewcommand{\thefootnote}{*)}
\newcommand{\fnmark}[2]{\renewcommand{\thefootnote}{#1}\footnote{#2}%
  \renewcommand{\thefootnote}{*)}}

% A piece's head, mirroring transcription.tex.  #1 = the English of the
% title as the volume's Indhold gives it (table of contents, running head);
% #2 = the head as printed, translated (a book under review keeps its Danish
% title, glossed in brackets).  The short rule under every head is printed.
\newcommand{\piecehead}[2]{\clearpage\phantomsection
  \addcontentsline{toc}{section}{#1}\markboth{#1}{#1}%
  \begin{center}\large #2\\[6pt]\rule{2cm}{0.4pt}\end{center}\par\medskip}
\newcommand{\piecerule}{\par\begin{center}\rule{1.5cm}{0.4pt}\end{center}\par}

\usepackage{hyperref}
\hypersetup{pdftitle={Brandes, Collected Writings XIII: The Controversy over Faith and Knowledge}, pdfauthor={Georg Brandes}, hidelinks}

\pagestyle{fancy}
\fancyhf{}
\fancyhead[LE,RO]{\thepage}
\fancyhead[RE]{\textit{Collected Writings XIII (1903)}}
\fancyhead[LO]{\textit{\leftmark}}
\setstretch{1.3}

\title{\textbf{The Controversy over Faith and Knowledge}\\[8pt]\large Georg Brandes, \textit{Samlede Skrifter} [Collected Writings], vol.~XIII (Supplementbind),\\ pp.~1--42}
\author{}
\date{Copenhagen: Gyldendal (F.\ Hegel \& Søn), 1903\\[10pt]
  \small Translated from the Danish}

% ---- original-page markers ----
% \opage{N} marks where printed page N begins, at the same joint as in
% transcription.tex (mid-sentence where the page turns mid-sentence).
\usepackage{marginnote}
\newcommand{\opage}[1]{\marginnote{\footnotesize\textit{[#1]}}}
\newcommand{\apage}[1]{\marginnote{\footnotesize\textit{[#1]}}}
\begin{document}
\maketitle
\thispagestyle{empty}
\tableofcontents

\mainmatter

% --- p. 1 ---
\piecehead{Introduction}{\opage{1}INTRODUCTION}
In the Preface to my \emph{Samlede Skrifter} it was remarked that the
edition would not include everything I had published in the course of my
life. Even so, it grew far more extensive than the original estimate had
promised, and included, among other things, besides pieces collected
earlier, 1200 pages that had not yet been published in book form.

Left out was everything polemical in form, everything that was not
immediately accessible to the reader, and a multitude of shorter
articles whose inclusion would have made the work swell far too much.

An author is badly placed, however, if he does not himself collect and
take care of whatever --- regardless of its greater or lesser worth ---
bears the stamp of his cast of mind. For, as experience shows, what he
himself lets slip out of his hands, other hands seize upon. Once a
writer has succeeded in making his name sufficiently known, strange
fingers rummage in his literary life and present it in whatever way
their presumed superiority sees fit. And this superiority has all the
freer play, the more he has left to its caprice.

This consideration has led me to publish a supplement to my Samlede
Skrifter, and I open this first supplementary
% --- p. 2 ---
\opage{2}volume with my earliest polemical articles, to be followed by a
selection of others, mixed with works that aim only to present and to
instruct. As regards these earliest articles, and the few barbs they
contain against certain distinguished authors such as Professors Høffding
and Wilckens, whose goodwill is an honor that I cannot be suspected of
wishing to forfeit, and whom I would therefore least of all wish to
offend, I may be permitted to recall the fine letter from J. L. Heiberg
to Carsten Hauch (28 January 1841), in which the former, after warmly
accepting Hauch's outstretched hand and assuring him that an old feud
between them has not left a spark of ill will in his mind, tells him that
a reprint of some polemical articles from that time ``can and should in
no way disturb the good relationship'' which, to his joy, now exists
between them: ``I know well that this renewal of what one might in one
respect wish forgotten is easily misunderstood''. But, he says, ``it
exists, after all, and in literature as well as in political history not
only can facts not be altered, but it is even a duty to collect and
preserve them.''
I do not go quite as far as J. L. Heiberg. Many polemical particulars are
best consigned to oblivion, and it would not occur to me, e.g., to
reprint my polemical sallies from the Nietzsche controversy of 1889--90,
now that Høffding has left his out of the collection \emph{Mindre Arbejder}.
But to some of the literary feuds in which I have been entangled people
will perhaps return in the future. Much that now seems to have only
historical interest may, when kindred phenomena arise again, become
topical once more.

Faced with a critical author, the reading public divides into two groups:
a larger one, which reads him if it is interested in the subjects he
treats, and a smaller one, which is interested in the author himself. In
the \emph{Samlede Skrifter}, as was expressly stated in the Preface, the
whole emphasis was laid on the subject matter.
% --- p. 3 ---
\opage{3}Regard for the general public's need for instruction was the
one deciding factor there. What belonged together by the nature of its
subject was never split apart, even if it had been produced at widely
separated times. And nothing was included that could lay claim to any
attention only by its manner of treatment.

In the Supplement, however, there should also be room for this and that
which can hold only those who already take an interest in the writer and
who read him even when the subject he treats is a matter of indifference
to them. The critical portrayal of a personality is, after all, an art in
its own right, and can be judged much as portrait painting is, without
regard to whether the reader or the viewer knows the subject or not. The
art of fencing in intellectual disputes is likewise a small art of its own,
and can amuse the reader as a fencing bout can entertain the spectator,
without regard to the prize for which the fencers fight.

In other words, besides the broader reading public that seeks information
about the subject matter, there is also a more finely developed public,
which is drawn more by the author's personality than by the subject on
which it happens for the moment to dwell. It is to this public that the
supplementary volume is addressed.

April 1903.
\begin{flushright}G. B.\end{flushright}
\piecerule

% --- p. 4 ---
% blank
\opage{4}

% --- p. 5 ---
\clearpage
\phantomsection
\addcontentsline{toc}{chapter}{The Controversy over Faith and Knowledge}
\begin{center}\opage{5}{\Large THE CONTROVERSY OVER FAITH AND KNOWLEDGE}\end{center}

% --- p. 6 ---
% blank
\opage{6}

% --- p. 7 ---
\piecehead{Conscience and Science}{\opage{7}\emph{A. C. Larsen: Samvittighed og Videnskab}\\[2pt]{\normalsize[Conscience and Science]}\\[2pt](1866)}
Writings of philosophical content are in general not especially suited to
review in a daily paper, and if some enthusiastic man now and then
ventures the attempt, it easily turns out, as experience shows, to be a
``scientific work of poetry'' in the style of Lieutenant v. Buddinge.
Since learned critiques are rightly regarded as inadmissible, the result
is that on scientific subjects our papers teem with extremely brief,
unproven and often enough unprovable assertions, which the public has
indeed acquired a true virtuosity in swallowing quite raw, but whose
number no sensible person can nevertheless feel any desire to increase.
There is thus not much to invite one to review a pamphlet that does not
belong to light literature; but since this one treats, with more than
ordinary ability, a matter that not only has a claim on wide interest but
has long since laid hold of it, the circle of readers the pamphlet may
expect ought to be made aware of its appearance.

It contains a fundamental theological outlook, set forth partly in the
form of communication, but more often by way of defense and attack, and
as a scientific confession it is of particular interest, since the author
possesses a philosophical cultivation that is becoming rarer and rarer
among the theologians of the present day, and that helps him past the
first reefs, on which the majority, with sublime composure, sit stuck
fast in a common shipwreck.

Although the book is at feud with living persons in only a few places, it
will nonetheless probably be regarded as a polemic, and Mr. \emph{Larsen}
will himself bear part of the blame, since in a preface he has laid
particular stress on his offensive
% --- p. 8 ---
\opage{8}aim. The part of the book that is of immediate interest is,
moreover, the best, and one must be grateful to the author for having had
the courage and the will to embark on a discussion of the boundary
questions concerning the mutual relation between faith and knowledge, on
which a clear and illuminating word is so seldom heard; among the
students, the need to have these questions answered has long since pushed
back every other need for philosophical learning, and not one of the
internal tasks of science lies half so close to their hearts as this
boundary dispute. Several, to be sure, have come to rest, but mostly
those who would feel entirely satisfied with any conciliatory solution
whatever; it is only human, after all, to want to have it both in the bag
and in the sack, and natural that one is then not too particular about
whether the bread in the bag is honest bread or not.
But far greater is the number of \emph{those} who, after the lapse of a
few years, have retained hardly any impressions of Professor
\emph{Nielsen's} brilliant and fruitful teaching other than those bearing
on the difficulties mentioned, and for whom these questions are still and
always burning. This state of affairs is neither good nor healthy, all the
more so as those impressions are never quite clear. For in Professor
Nielsen's philosophy, which still unfolds luxuriantly before the students'
eyes, this point is, in the strict sense, the salient point, the real
``unrest'' that set the movement of thought going, drove the problems
forward and kept the gifted worker from seeking rest for a single moment
in rigidified forms.
But it has seemed as though it were a condition of the life of the whole
that this innermost provoking question should be kept constantly open,
that the wound should not close where the thinker had been most deeply
wounded by the god. Not that no answer was given; but the answer is
better suited to calling forth new questions than to silencing those
already raised.
% "have baade i Pose og i Sæk" = to have it both ways; "have rent Brød i
% Posen" = to have clean (honest) bread in the bag, i.e. to have a clear
% conscience.  Brandes joins the two idioms; the English keeps both images.
% ‹Uro› is both "unrest" and the balance wheel of a watch, which keeps the
% mechanism going; hence "gave the movement ... Fart".  The English cannot
% carry the second sense.

Already in the Professor's \emph{Propædeutik} of 1857 the sections that
treat the relation in question contain a plain self-contradiction: they
let faith and knowledge, though ``absolutely heterogeneous principles'',
be reconciled in one consciousness, since the conflict between them
``does not arise from what pertains to principle'', and this
contradiction has not since been removed. Professor Nielsen's position
here is altogether such that he can expect to find two kinds of opponents:
first, those who defend the existing theology, and
% --- p. 9 ---
\opage{9}secondly, those who attack it at its root and ground. The former
will go after Prof. Nielsen as the assailant of theology; the latter,
who seek to overthrow theology, must strive to establish their right to
call Prof. Nielsen himself a theologian. The opponent who has now
appeared stands in the ranks of the former. He fights as a theologian for
hearth and altar, without rancor or malice, but not without an eye for
the openings his superior adversary presents. He does not succeed in
proving what he had actually set as his aim, namely that theology still
has vitality, and there may well be good reasons why such a proof was
bound to fail; but he establishes beyond doubt that Prof. Nielsen cannot,
from his standpoint, reject theology, even if a certain particular
theology and a certain particular dogmatics have been exploded.

The author has one further merit. The natural reverence for Kierkegaard's
memory has imposed silence about his works; they have been allowed to
speak without any buzzing swarm of reasoners disturbing the impression of
the mighty monologue. That this was right, every Dane will understand; a
German would never grasp it, and Kierkegaard's Norwegian trumpeter,
Mr. Dietrichson, perhaps least of all. But all the same, it was the right
thing. Yet a drawback has followed from it, namely a petrifaction of
certain Kierkegaardian propositions in seemingly trustworthy
interpretations. On such points, then, an inquiry into Kierkegaard's
meaning must be set going, and in the essay on the Paradox Mr. Larsen has
begun one.

But once these merits have been brought out, the best there is to say
about the book has also been said. The author, in whom, despite a
predominant independence, one can trace the influence of Jacobi and
Schleiermacher, is lacking still, to some degree, in schooling, and to a
high degree in dialectic. The first lack betrays itself, e.g., when at the
end of his essay on conscience he demonstrates the impossibility of
giving an objective proof of God's existence so thoroughly that in the
same stroke he ends up demonstrating the impossibility of inferring from
effect to cause at all; or when, in his somewhat layman-like feud with the
idealists, he proves the existence of lifeless objects from the existence
of the self-conscious being before this being has become conscious of
itself. The lack shows itself above all in a sentence like this: ``The
miracle precisely confirms the law of nature; if one posits the exception,
one must also posit the rule'', in which the expressions \emph{law} and
\emph{rule},
% --- p. 10 ---
\opage{10}just like the two famous sharks, swallow each other up to their
hearts' content.

Of still greater significance is the lack of dialectic. The author speaks
the language of philosophy, he uses its designations, he romps about with
its expressions; but --- the voice is not its voice. In every line one
senses the theologian.
The author's ``Absolute'' is --- at times astonishingly --- conceived
undialectically, and why? Because in his mouth this expression does not
designate the Idea of philosophy, but is merely a name, a hiding place, a
mask behind which he seeks to make theology's representation of God
unrecognizable. But since the designations of concepts react upon the
train of thought, he not infrequently entangles himself in contradictions;
at one moment, e.g., man is supposed not to differ in essence from the
Absolute (the Eternal, the True), at another he is supposed not even to be
able to comprehend the possibility of the Absolute's actions. When
explanation now stands against explanation, philosophy's against
theology's, the author takes up the following position. Since he is short
of dialectical ability, the philosophical explanation becomes
unintelligible to him; the miracle explanation, of course, does not even
claim to be intelligible; and so it comes down to a choice between them.
For this choice he then holds not reason but the heart responsible. That
one will not accept the miracle explanation does not confuse him, then;
that he understands perfectly well, for it rests on the hardening of
hearts. Here Mr. Larsen casts aside his philosophical garb and shows
himself in his theological vestments. Here, then, is also the limit of
scientific discussion.
% "men --- Røsten er ikke dens": probably an echo of Gen. 27:22 ("The voice
% is Jacob's voice, but the hands are the hands of Esau").

December 1865.

\emph{Postscript}. In an article on Mr. Larsen's book in yesterday's
\emph{Dagbladet}, Mr. \emph{Rudolf Schmidt} declares that it is not his
intention ``on this occasion'' to clear up the obscurity in Professor
Nielsen's doctrine mentioned above. This sad news is hereby conveyed to
the readers of \emph{Fædrelandet} as well.
\piecerule

% --- p. 11 ---
\piecehead{Reply to Mr. Schmidt}{\opage{11}\emph{Reply to Mr. Schmidt}\\[2pt](1866)}
Besides an unsuccessful defense, Mr. Larsen's pamphlet contained an
attack that carried the war into the adversary's own country. The cut was
incomparably better than the parry; for a good cause arms eye and hand.
In my Saturday article I endeavored to set the justification of the
attack in a sharper light, and I held fast to this point as the main
point; for it is on this that the interest rests. The battle against
theology is of old date; \emph{Feuerbach} and \emph{Kierkegaard}, each
from his own side, waged it with a force and consistency that have never
since been attained; little weight now rests on that battle; but as
surely as Professor \emph{Nielsen's} philosophy is a far more significant
and far more intelligent power than the theological science of our day,
so surely does every doubt concerning it that is not plucked out of the
air, and every thrust that does not strike the air, have a great and
essential significance. But the matter deserves still greater attention
when its point is a center, when the question is a vital question of the
utmost importance and of pervasive influence; and this matter is of that
kind. It is this that Mr. \emph{Rudolf Schmidt} treats as a trifle.

Now if the contradiction I have pointed out had really been removed in
any of Professor Nielsen's printed writings, then I would certainly be a
greater charlatan in philosophy than any of the several that our
fatherland possesses. But that is not how matters stand. Mr. Schmidt is
kind enough to pay me a compliment on my philosophical cultivation, one
that has a reference to \emph{Grundideernes Logik} in tow. I have as
little use for the latter as for the
% --- p. 12 ---
\opage{12}former; I do not even understand what Mr. Schmidt had in mind
with it. In the published part of \emph{Grundideernes Logik} there is,
apart from a few propositions put forward as bare assertions in the
general introduction, nothing, absolutely nothing bearing on this
matter, nor, given the nature of the design, could anything be found
there. Only at the transition from the sphere of knowledge to that of
power can the question come up, if in a form more remote from reality,
and it can receive its proper answer only at the transition of the whole
Logic into the philosophy of nature.

I did not expect enlightenment from Mr. Schmidt; I merely made fun of an
expression of his. I have a well-founded doubt that a polemic with him
would prove instructive, and shall therefore not engage in one; for what
I want is to learn. But for that very reason an explanation on
Prof. Nielsen's part would be of the highest value to me, and nothing
would please me more than if he felt called upon to give one. For in
this matter what we need is not silence or reticence, but words that
come from the liver and go to the heart, unsparing candor, clarity,
light:
% "Ord, der gaar fra Leveren": "tale fra Leveren" (to speak from the liver)
% is the Danish idiom for speaking frankly; kept literally for the pairing
% with "to the heart".

\emph{Truth, like a torch, the more it's shook, it shines.}
\piecerule

% --- p. 13 ---
\piecehead{Philosophy and Theology}{\opage{13}\emph{On Mr. Harald Høffding's}\\
\emph{Polemic: Filosofi og Teologi}\\[2pt]{\normalsize[Philosophy and Theology]}\\[2pt](1866)}
When one reads through this work to make a passing acquaintance
with the author's features, one receives an almost uncanny
impression; for with single and double quotation marks (``goose-eyes,'' as Danish calls them)
% PUN: Danish `Gaaseøjne' (goose-eyes) is the ordinary word for quotation
% marks, whence the ``Argus of the goose tribe''; the gloss is added so the
% image can stand.
by the hundred --- sometimes as many as thirty on a
single page --- the author, like an Argus of the goose tribe, stares
gloomily out at the reader from the paper of the book. With him, that
is, the quotations play the same role as the text of one's own making
plays with other authors.

Proceeding from the assumption that the philosophical-theological
controversy has in reality already been settled, Mr. \emph{Høffding}
is concerned only to make it clear that Prof. \emph{Nielsen} is the
one who has won; and it is to this end that he gives a heavy,
densely packed restatement of the history of the battle, which, % Danish prints `forkjellige' for `forskellige'; rendered as the sense.
according to the reader's differing standpoint, is either insufficient or
unnecessary. Against Mr. \emph{Larsen's} pamphlet
% The Danish opens this quotation with a turned (closing) mark; the English
% balances.
``Samvittighed og Videnskab'' he then advances
a number of arguments, mostly ones that have already been
seen in print once before, and finally directs against the undersigned's
earlier articles in \emph{Fædrelandet} some objections which make up the
weakest, but perhaps also the newest, part of the book. For these
objections I am very grateful to the author; for they fully demonstrate
what I have said before, namely how badly an explanation is needed
from Prof. Nielsen's side.

I have declared it a contradiction that Prof. Nielsen
at one and the same time calls faith and knowledge absolutely heterogeneous principles
% --- p. 14 ---
\opage{14}and speaks of their reconciliation in human
consciousness. Mr. Høffding teaches me that in Prof. Nielsen's view
this reconciliation can take place, not \emph{although}, but \emph{because}
these principles are absolutely heterogeneous. This is rather far from
being any secret; but the contradiction is designated all the same by
my ``although.'' Is it conceivable that two absolutely heterogeneous
principles can be reconciled in the same consciousness? Can it be
shown that faith and knowledge are such principles? I answer these
questions in the negative. It might, however, well be open to doubt
whether two absolutely heterogeneous principles could hold sway in the
human spirit without a split setting in by which the whole of existence
would be cracked apart; and that Prof. Nielsen himself has not taken
this absolute heterogeneity so very strictly is best proved, after all,
by the circumstance that he speaks of conflict and reconciliation between
those powers; for absolutely heterogeneous powers can neither wage
war nor make peace with one another. But let us for once
suppose that we have two absolutely heterogeneous principles before us:
does Mr. Høffding not see, then, how strange all this talk becomes
about how one is to place faith highest and science lowest,
but not the reverse? Does he not see that every question of
what is to be placed highest and what lowest would in that case become
still more comical than \emph{Poul Møller's} old one: Which is higher, a
poplar-willow or a thunderclap? And yet many a poor student has perhaps
come to grief when, at the examination drill, he answered
wrongly and in his innocence turned faith and
knowledge upside down.

To ease understanding, however, Mr. Høffding gives
yet another hint. Ethics, philosophical ethics, ``forms a transition
between philosophy and religion.'' Considered in and for itself,
this proposition may be quite excellent; but in Mr. Høffding's
mouth it certainly cuts an incomparable figure. For morality has no doubt
always been regarded as an admirable power; but that it could form
the transition between two absolutely heterogeneous spheres---that
is something hardly anyone has credited it with. Rather than
investigate how this was supposed to come about, one is tempted to
ask what moved Mr. Høffding to send this piece of cleverness out into
the world. The answer is not difficult.

To anyone who has followed with an attentive eye the course of
Professor \emph{Nielsen's} rich activity as an author over the last decade,
it has gradually become clear that, without taking back a single
% --- p. 15 ---
\opage{15}spoken word, he has, on the disputed point, if not changed, then at
least altered and softened his earlier view. I am far from imagining
that I am saying anything disparaging by this. An avowal that one has
changed one's conviction from the ground up can be a noble, indeed an
admirable, trait in an eminent personality, and such this act proved
to be when Professor Nielsen in his time attached himself to
Kierkegaard. To develop the same fundamental thought while adapting it
is, if anything is, not merely an honest but an honorable
thing. The latter, in my understanding, is what Professor Nielsen has done.
In 1857 he declared that between the evangelical and the mythical
there was no difference in the eyes of science; but in 1864 he softened
this expression so that the superiority of religious ethics over the
rational prevented science from declaring the Gospel
a myth. The latter is indeed a hard saying; it is difficult
to understand that there can be two kinds of morality, and that
any other morality can be higher than the rational, i.e. the
morality grounded in the rational cognition of man's nature. But no
matter: it is plain that a softening has taken place here; for with this
the absolute heterogeneity of faith and knowledge is at bottom given
up, though not retracted.

For such things Mr. Høffding has no eye. A development which,
despite the apparent protest of particular utterances, tends in the
direction of raising science from the subordinate relation in which it
was formerly placed---this he cannot see; instead he boldly heaps
propositions torn loose from the most various books into a pile and
stacks one contradiction on top of another. He lays it down that neither
the right, nor the true, nor the good could hold its ground in real life
if it were not supported by ``religion''; but he is nonetheless
convinced that there is a philosophical, hence rational, morality,
which --- on the assumption just mentioned --- will at least have the less
rational property that one cannot live by it. He maintains with the
utmost zeal that the demands of conscience must not be asserted in
science; but nonetheless, with the greatest unconcern, he reproaches the
freethinkers' science with being an ``unbelieving'' science. Mr. Høffding
is no friend of consistency.

On the last page of the work the author raises the objection
against me that, if I have found no contributions to an answer to my
question in what has so far appeared of \emph{Grundideernes Logik}, then I ought
% --- p. 16 ---
\opage{16}not to have confined myself to wishing for enlightenment from
Professor Nielsen, but should have found fault with what was given. This
objection was easy to foresee and is easy to refute. It is, to be sure,
incomprehensible to me how Professor Nielsen can let the possibility of
miracle obtain a place in his system without running aground in final
disunity; but in the part of the \emph{Logik} that has appeared, no attempt
whatsoever has been made at that either. Nor do I see
% The Danish prints a point after `bebuder' where the sentence runs on;
% rendered as the sense.
how the development is to be able to give what the preface announces
without the higher unity of knowledge and power, of ideality and
reality, which is the fundamental thought of the work, being burst
apart---most immediately by the idea of power turning into the
representation of omnipotence in the theological sense. I alluded to
this in an earlier article.

There are in this country, as elsewhere too, e.g. in
France, certain authors who can best be likened to those well-known
large bladders which are in themselves flat, slack and empty
things, but which, with the help of inspiration from outside, swell up
so that they take up quite a considerable space. It should be said
% The Danish prints `Heffdings' (wrong-fount e); rendered as the name.
to Mr. Høffding's credit that he does not belong to this class. He
evidently has an honest endeavor and a serious mind. But
these merits have not prevented him, in his relation to Professor
Nielsen, from showing himself on this occasion to be an all too
loose-tongued child.
\piecerule

% --- p. 17 ---
\piecehead{A Last Rejoinder}{\opage{17}\emph{A Last Rejoinder}\\[2pt](1866)}
Every theatergoer will have had the experience that even the
bad actor can reap claps and applause when, as the knight of
innocence or of faith, he pronounces a noble principle,
and that the audience sometimes listens almost with displeasure as the actor
defends a cause it considers bad. If my honored opponent,
Mr. \emph{Høffding}, could get the roles so distributed between
% The Danish prints `at kan kom' for `at han kom'; rendered as the sense.
us that he came to play an angel of light and I a demon of
darkness, then his cause would undeniably have won, and outward
triumph would have been assured him. Such a turn he tries to
give the matter by accusing me of not wanting to be a
party but a judge, of avowing a freethinking view not openly and
honestly but only by roundabout ways, when after all ``no one
may take part in the contest without declaring what he bears on his shield.''

But this turn shall not succeed for him; it is not unbelief
against faith, rather it is clarity against fog, that here stand face
to face. For what is the object of the inquiry, what alone is
being asked in this matter? Certainly not about \emph{my}
philosophy, any more than, e.g., about Mr. Høffding's. There is only one
single decisive point here which must be held fast, and
it is this: Is Professor \emph{Nielsen} in a contradiction when he
calls faith and knowledge absolutely heterogeneous principles and yet lets
them be reconciled in human consciousness, so that the
one is placed highest, the other lowest, or is he not? He
is; for either faith and knowledge are not absolutely heterogeneous, or
they split consciousness completely; either they are not absolutely
heterogeneous, or there can be no talk of a highest and a lowest. It
% --- p. 18 ---
\opage{18}is no use objecting, as Mr. Høffding does, that it is after all the
same human being who stands in a relation to both of the two absolutely
heterogeneous principles; for the question is precisely whether
man's consciousness must not thereby be incurably divided. Nor does it
help to say, as Mr. Høffding does, that for this very reason the
question is put to each person individually how he will make
his choice here. For what does the choice consist in? If the choice is between
faith and knowledge on the one side and knowledge on the other,
then whoever seizes the former chooses, after all, the split; if on the
other hand the choice is between faith on the one side and knowledge on
the other, then reconciliation in the same consciousness is, after all,
given up, and only dualism, only disintegration, remains, even if it
is shifted from the individual human being to mankind. It is
of no use at all, like Per Degn, to answer simple and clear
questions with highly enigmatic and unintelligible talk.\fnmark{*)}{So that Mr. Høffding shall not once more send my own turns of phrase
back to me within the marks I will not name, and see it proved here into
the bargain that I treat Professor Nielsen as a ``delinquent,'' I add that from this
comparison it follows just as little that I am on the whole any Erasmus
Montanus as that I have, like him, perpetrated the outrage ``of striking Rasmus
% The Danish never closes the quotation opened at `at slaa'; the English
% closes it after the inner one.
Nielsen 3 or 4 times on the back of the neck with clenched fists, all the while shouting:
`\emph{Probe Majoren, probe Majoren!}\,'\,'' (Act 1, Scene 6).}

I do not consider myself infallible, as Mr. Høffding gives his
readers to understand; I do not even consider Professor
Nielsen so, which Mr. Høffding would doubtless have forgiven me
far more readily; but if I am not so immodest as to
take it to be impossible that I can err, I am still far
less so foolish that I should have set pen to paper, without a rock-firm
confidence in the truth of my cause, against a man so highly gifted
and so superior to me as Professor Nielsen.
What I bear on my shield, then, is this: a conviction,
full and firm, that every dualism (twofoldness) is impossible,
and that in consequence Professor Nielsen's doctrine on this
point is untenable. What philosophy I have besides, or whether
I have any at all, is quite as irrelevant to this question
as what I think about the affairs of the
Turks.

I have strengthened my argument by showing that Professor
Nielsen himself has not been able to hold fast to the absolute
heterogeneity. If it really obtained, then any
% --- p. 19 ---
\opage{19}form of religious faith whatsoever would have to be compatible with
Professor Nielsen's philosophy. But if one asks, e.g.: can one at the same time
be a Catholic and an adherent of Professor Nielsen's doctrine? then
the answer is Yes only on the condition that faith and knowledge have
nothing whatever to do with each other, but No if account is to be taken of
the ethics of Catholicism; for it is surely beyond doubt that this
is not superior to the rational. In making the superiority of the ethics
of religion the decisive thing, Professor Nielsen has given up his
first standpoint. Since Mr. Høffding, however, believes that he
can defend the proposition of the double ethics as a transition
between philosophy and religion, I shall gladly follow the curious
course of his thoughts.

There is, he says, an opposition between science and
life. When the idea of the good is to be realized in deed,
this opposition comes forward in such a way that rational ethics shows
its insufficiency, and the necessity is seen --- not of religious
ethics, as one would provisionally have expected, but --- of
``religion.'' Now I ask: what, then, is meant by religious
ethics? Is it a science, or is it merely another name
for the positive-religious? If it is the latter, why then the
confusing designation? But if it is the former, then the matter is
simple and the confusion plain. For then, by virtue of that very
same opposition between science and life, this ethics too must
become insufficient as soon as it is to be used in life, and
recourse must again be had to something higher. But, not to mention that
the significance of rational ethics, which was after all to be upheld,
is by this line of thought reduced to less than nothing, we get
here a series of the scientific, the superscientific, the super-super-scientific,
etc., which does not speak in favor of the logic of this
defense. If being a moral human being involved such
logical entanglements, then I do see that one would have to ask with Mr.
Høffding: Can one at the same time be an ethical man and a
man of science? But otherwise this question is not easy to
understand. --- No, it stands firm: if ethics makes the transition from
philosophy to religion, then their absolute heterogeneity is given up, and
the best proof of this doctrine's untenability is that
Professor Nielsen has given it up not only here but, as
I have shown, on other points as well. Thus,
when the possibility of miracle is given a place in \emph{Grundideernes}
% --- p. 20 ---
\opage{20}\emph{Logik}.\fnmark{*)}{I was first informed that the solution of the problem
was given in the part of the \emph{Logik} that had appeared. I replied that
it was not to be found there and, given the plan of the work, could not be
found there. The same people who picture faith and knowledge shut off
in two different compartments of consciousness then accused me of
subscribing to a pigeonhole philosophy. I replied that I understood the
fundamental thought of the work to mean that the possibility of miracle
could not be prepared for or derived, but could come about only by a leap.
Now Mr. Høffding does not refrain from reproaching me for not having
pointed out the error in what has already appeared---an error which I myself
declare is not to be found there.}
For the difference between the Absolute of philosophy and the
almighty God of theology is this, that the latter is not, like the Absolute,
conditioned by its own conditions; but this difference the Logic
abolishes in positing the possibility of miracle. In the Logic,
then, the same doubling appears as in the ethics. Just
as it cannot be understood how religious ethics
can fail to have a retroactive and annihilating effect
on the rational, so it cannot be conceived that the
unconditioned personality, by virtue of its self-determination, can raise
itself above its own self-determinations without their
significance being annihilated, that is to say: without the laws of nature
and of spirit, as we know them, becoming nothing, mere semblance. I
understand that one can, with Kierkegaard, support one's paradoxical
faith by rejecting science; but I do not comprehend how one can,
like Professor Nielsen, give up doubt about science and
yet keep that faith---at least not without being compelled
to assert against the freethinker those demands of conscience
which one will not tolerate the theologian's directing at oneself. This
statement may perhaps explain to Mr. Høffding what he is so
very interested in knowing, namely what it is that I
recognize in Kierkegaard. Only people of very little independence, surely,
understand by recognition nothing but a plain and simple
adherence through thick and thin. What I admire in
Kierkegaard is, among much else, his consistency with
himself.

This is my last word to Mr. Høffding. Now nothing, to be sure,
is further from my mind than to conclude from Professor Nielsen's silence that
he will endorse what my opponents have asserted
against me, or that their contributions really contain
the correct interpretation of his view; but I believe
that it would have served the cause if he had
% --- p. 21 ---
\opage{21}made this known. A tangible result a contest of ideas like
this one could not bring; but it is my hope that, if my statements
have yielded nothing else, they have at least perhaps made one
or another who strives in quiet---one who has the courage to judge, the will
to judge honestly, the understanding to judge rightly---aware
of where the difficulty lies.
\piecerule

% --- p. 22 ---
\piecehead{On Self-Responsibility}{\opage{22}\emph{On Self-Responsibility}\\[2pt](1866)}
The author of ``Julen hos Elise'' has brought out, with Gad's publishing
house, a work, ``\emph{Om Selvansvarlighed}'', which may serve as an
example to illustrate the following sentence in ``\emph{Enten---Eller}'':
``There is a kind of reasoning chatter which in its endlessness stands in
the same relation to the result as the interminable lines of Egyptian
kings to the historical yield.'' For these words give the right point of
view for a book which, like this one, is without force in its expression
and without movement in its thoughts, boneless in its structure, jointless
in its articulation and bloodless in its execution.

The author has set himself the laudable aim of being generally
comprehensible, but he has not seen that comprehensibility is attained far
better by a discipline of thought that excludes foggy breadth, irrelevant
digressions, repetitions and triviality than by any of these scarcely
brilliant vices themselves; he has further followed the laudable plan of
leaving all earlier works on this subject unmentioned; but it almost seems
as though he had also himself either not read them or not understood them;
for otherwise it is incomprehensible that he has not even been able to
\emph{pose} his problem with any precision, even if he could neither solve
it nor cut it through. His work, then, belongs as little among those that
make the results of a science clear to the great multitude as among those
that carry literature forward; it ought not to have been published as an
independent book; in abridged form it could have claimed a place in one of
our two theological journals.

The fundamental mood of the book is a deep sense of the inadequacy of our
faculties to solve the riddles of existence. The
% --- p. 23 ---
\opage{23}author always begins by theologically shrouding matters in mystery,
only to declare them inexplicable afterwards. Of \emph{conscience}, for
example, it is said that although it must stand in a continuous connection
with the all-governing principle, it is ``\emph{entirely inexplicable} how
% The Danish closes this quotation with a turned mark (`vedligeholdes‹');
% the English closes it normally.
the connection takes place and is maintained''; of \emph{life} we are told
that the vital force (!) makes itself known only when it is joined (!) with
a bodily existence, ``but how the vital force joins itself with the bodily
in such a way that this can reproduce itself in a new individual is the
\emph{Unfathomable}.'' Sometimes he presupposes degrees of intelligibility
within this Unfathomable, as when in the following sentence he arbitrates
between theism and pantheism: ``If the first origin and the entire activity
of a living, rational, personal being as almighty God is
\emph{incomprehensible}, then the origin and productive power of an
impersonal fundamental force is no less so; but the life of the world, its
purpose and goal, becomes \emph{nevertheless more intelligible} when it is
regarded as proceeding from a personal God.'' Sometimes, on the other hand,
he gives up the struggle altogether, as in the following despondent, naive
declaration, though a natural one in the author's mouth: ``we confess that
we can never attain any clarity in our thought about the things that cannot
be presented to our senses.''

To gain a foothold in this darkness the author now sometimes starts from
theological presuppositions, such as that of an original state of morality
which was later lost, or that of original sin; sometimes he appeals, as his
first presupposition, to our \emph{feeling:} positive religion is confirmed
by ``our own inner feeling'' or by ``the involuntary feelings of a purer
human nature,'' and an opposing view, which admittedly ``may well have
something in its favor in what we experience by way of the senses about the
course of development the world has passed through'' [that is, in the one
thing about which, in the author's opinion, we can attain clarity], is
refuted by the fact that ``there is nothing in it that can appeal to human
feeling.'' Once this criterion of truth has been adopted, woe then to the
doctrine of infusoria and trichinae; for truly these creatures have nothing
about them that appeals to human feeling.

On the foundation now indicated, the relation between faith and knowledge
is determined as follows: ``The one is won by a practical path, through
experience, by consideration of the visible, which can be observed by all;
therefore the representation of it can also be
% --- p. 24 ---
\opage{24}compared with, and corrected by, the representation others have
formed of it, and attain such certainty that we can designate it as
\emph{knowledge}. The representation of the higher (!) spiritual, on the
other hand, becomes of a more (!) subjective nature, because it has nothing
visible to attach itself to that is observable by all in the same way; we
designate it as \emph{faith}.'' By this fine and precise definition of
concepts, logic and thoroughbass, among other things, become faith; but
with all the amiable unselfconsciousness of a young modern lady the author
is sure of being understood when he hangs faith upon the capacious
comparative that goes by the name of ``the higher.''

The author is as great a stylist as he is a philosopher. In the first
capacity we see him when he teaches that ``many passions and affects can
arise and persist only with \emph{bodily support}, such as irascibility and
a propensity to debauchery'' (it takes a body), or when he speaks of how
the human body ``expands from a most insignificant beginning constantly in
greater and greater compass.'' As a thinker we hear him in the following
priceless derivation of the love of honor: ``The beautiful body arouses
admiration, but only the human body can call forth honor; for it is the
spirit alone that, as the possessor of this body, becomes the object of it.
It is chiefly the spiritual expression, which makes itself known most
immediately in the features of the face, that can call forth honor. But the
expression is uncertain and cannot in everyone penetrate the corporeality;
therefore [ergo] man must seek to win honor by his actions.'' Holberg's
Henrik never conducted a proof more wittily.

In the book's seventh section, ``The Relation between Human Freedom and
God's Governance,'' the author has concentrated the results of his research
% The Danish prints no full stop after `Forvirring'; the English supplies it.
and his whole confusion. This section justifies the Deity. The reasoning is
as follows: The governance that reaches its goal while preserving the
freedom of the persons it uses as means must be regarded as the most
perfect. Consequently man's freedom does not conflict with the perfection
of the governance. Since, however, a development proceeding by virtue of
freedom cannot in every particular accord with the divine will, the
question arises whether it would not have been most worthy of the
governance to cancel the evil effects of freedom, often calamitous for the
innocent, by an extraordinary intervention, by the miracle. The author
denies this. Why?
% --- p. 25 ---
\opage{25}His manner of reasoning is worth hearing: ``If the shot aimed at
the innocent man's breast glanced off it without harming him, the result
could stand as a testimony to his innocence, but it would not cast any
light on the act committed; the act could be set up as a test of the nature
of the one against whom it was directed; it would not serve toward a better
cognition of the good or the bad in the intention that led to the act.''
How such a cognition of the moral character of the intention is to be
derived from the outcome, we leave it to the author to prove; and merely to
be sure we understand him, we put this question to him: Is there,
then, no miracle, nothing supernatural?

One must surely be astonished that a theological philosopher as well
brought up as the author seems ready to answer with a flat denial. But let
the reader be reassured: he does not do that either; in philosophy he forms
himself on the best models. He holds with Salomon Goldkalb ``vieles for und
% Salomon Goldkalb, the Jewish merchant of Heiberg's vaudeville `Kong Salomon
% og Jørgen Hattemager', speaks a German-Danish jargon: roughly "much before
% and behind, and yes and no at once".  Left as printed.
bag, und ja und nej tillige'': ``If we see no reason to assume such
breaches of the natural order and law as necessary for the completion of
God's plan for the world, this nevertheless does not exclude a divine
activity that must stand as supernatural for man. Such there must be; for
without it the power laid down in creation could not be maintained (!).
God's activity as sustainer of the world is supernatural in so far as it is
\emph{incomprehensible} to the human faculties, but it becomes no less
natural when it expresses itself, while holding fast to definite laws,
through a \emph{higher} nature.'' Crafty author! There at a single stroke
he found a use for his \emph{Incomprehensible} and his \emph{Higher}, and
indeed got the two fused together into a higher
incomprehensibility.

As one among many testimonies to the theological-philosophical quagmire in
which we are stuck here in Denmark at present, this book is not wholly
without interest. But for the rest we know no better advice to offer its
author than that contained in this old couplet from \emph{Peder Paars}:

\begin{center}\small
Your finger in the times you must stick, and smell\\ % I Fingren stikke maa i Tiden, lugte til,
what time you're living in, if you would heed me well. % hvad Tid I lever i, om mig I lyde vil.
\end{center}
\piecerule

% --- p. 26 ---
\piecehead{The Validity of the Idea of Nationality}{\opage{26}\emph{Claudius Wilckens: Nationalitetsideens Gyldighed}\\[2pt]{\normalsize[The Validity of the Idea of Nationality]}\\[2pt](1867)}
Against a treatise by Dean Kofoed-Hansen, \emph{Nationalitet og
Kristendom}, the author stresses that the question of the justification of
the idea of nationality vis-à-vis Christianity is only a single side of the
question of the justification of human culture as such, so that one cannot
deny that idea significance and validity without passing sentence on the
whole of culture. To that extent Mr. Wilckens may well be entirely right.

But when he now goes on to set forth his own conception of the relation
between Christianity and humanity, his view is even more unclear and
self-contradictory than his opponent's. From the standpoint of philosophy
he proclaims --- seemingly without properly knowing the meaning of the word
he uses --- an unconditional \emph{dualism} of divine history and human
history, of religious development and human development, of Providence and
world-governance. But dualism (fundamental cleavage) is the private shame
of scientific thinking, which the thinker himself ought surely to be the
last to put on display. Although the author himself declares that eternal
Providence \emph{set} the world-governance to rule over the destiny of the
nations, he stresses in spaced type that world history, guided by the
world-governance, has \emph{nothing, absolutely nothing} to do with
Providence, and Providence \emph{nothing, absolutely nothing} to do with it.

This is already going very far; but the author draws the bow tighter still
when, reproaching Dean Kofoed-Hansen for lacking even the faintest inkling
of the said remarkable dualism, he accuses him of being in the grip of the
same fundamental error that in the Middle Ages brought mankind
% --- p. 27 ---
\opage{27}to establish \emph{duæ veritates}. For who has two truths, if not
Mr. Wilckens?

In this doctrine of the relation of truth the author attaches himself very
closely to Professor R. Nielsen. When he occasionally designates his view
(thus, e.g., the curious proposition that ``God set the laws of nature to
bear witness against himself'') as an ``individual explanation,'' the
meaning of this is perhaps that, though not in such particular points, yet
in what he otherwise puts forward, he regards himself as the organ of
Professor Nielsen's teachings. If he is right in this, then what Holberg
says in \emph{Niels Klim:} applies to Denmark now: ``There are a number of
the learned men of Europe who cultivate both theology and philosophy with
equal zeal; as philosophers they doubt all things, as theologians, on the
other hand, they dare deny nothing.''
\piecerule

% --- p. 28 ---
\piecehead{Introduction to Rational Ethics}{\opage{28}\emph{P. S. V. Heegaard: Indledning til den rationelle Etik}\\[2pt]{\normalsize[Introduction to Rational Ethics]}\\[2pt](1867)}
In the wider sense of the word \emph{introduction}, any preparatory course
through which an author leads his readers into a science may be called an
introduction to it. Both a history of ethics and a history of morality
could, e.g., provide a preparation for the study of the science of the
good. But of such propaedeutics there is, for any doctrine whatever, only
one that is indispensable, namely the one that has the doctrine itself as
its object. The introduction to ethics that is not arbitrarily chosen but
necessarily required answers the question: What is ethics? The proper
problem of ethics itself is: What is the good? Since the author,
immediately after establishing the concept of ethics, proceeds to develop
the essence of the good, he has given no introduction in the narrower and
proper sense.

But, one might say, never mind the name: what does the book offer? To this
the answer is that Dr. Heegaard, who takes no account of chronological
order, has not wished to give a history of ethics; he has wished only to
use particulars of history to illustrate the various possible ethical
principles. To derive these is the task of his work. Since the principle of
division that is laid down, namely the opposition between the objective
and the subjective, is, however, an opposition that shows itself in every
idea and not merely in the ethical one, the good is not considered
according to its peculiar essence, but only from its general logical side;
the specifically moral concepts, such as duty, virtue, conscience, serve by
the plan only to illustrate the said opposition, so their ethical
peculiarity does not come out, and
% --- p. 29 ---
\opage{29}like the fundamental opposition, so too the one between actuality
and essentiality and the one between immediacy and reflection, on which the
subdivision rests, are purely metaphysical.

But, one might say, here after all is a principle of division; if only
\emph{that} is held fast, one ought not to take the author further to task.
Unfortunately it is dropped as quickly as it is set up. The ethics of
essentiality is not subdivided, subjective ethics is not divided at all,
and yet there was absolutely nothing to prevent the chosen points of view
from being applied. Even if there were not historical representatives for
all the forms --- which, however, can easily be found --- the cognition of
principles, on which the author lays so much weight, required that they at
least all be set out.

Still, Dr. Heegaard perhaps wishes the reader to concede that it matters
less whether the principle of subdivision changes, provided only the
principle of the main division is held fast. We therefore turn our
attention to this. But it is surely clear that if subjective-objective
ethics is to be the third member in the series, it cannot have forms under
it that are again declared to be one-sidedly subjective, like Kierkegaard's
ethics. That it nonetheless has them is intelligible only from the fact
that the author has here changed his principle of division altogether. The
historical phenomena presented in the third section are found together not
because their ethics bears a subjective-objective stamp, which on the whole
it does not, but because, in contrast to the preceding ones, they take up a
decisive relation to the Christian-religious.

This is not felicitous, but let it be! We will follow the author without
dispute here too, and merely test what forms must arise when both the
rational and the positive-religious are given as presuppositions. One can
then --- to use the author's expression: faith and knowledge --- either
reject faith for the sake of knowledge (like Feuerbach) or, conversely,
knowledge in favor of faith (like Kierkegaard), or one can assert the
validity of both, and then either in unity (like theology) or in
unconditional separation (like Nielsen). Of these forms Dr. Heegaard skips
theological ethics entirely: an ethicist such as Schleiermacher is not
included, whereas Nielsen, who has written no ethics at all, is treated at
great length. Should the author now say that theological ethics is too
% --- p. 30 ---
\opage{30}insignificant to be included, this will be no reason in his mouth,
since as the representative of one of the forms he has set up Feuerbach,
of whom he himself remarks that he has no ethical principle. Should he say
that theological ethics is a form that no rational person need take into
account, he must remember that the question is still undecided whether the
same cannot be said of the tendency that would hold faith and knowledge in
absolute heterogeneity. In any case theological ethics cannot be rejected
except through a critique; to shut the door on it is a way out, but not a
scientific one.

The new principle of division was the opposition between the rational and
the positive-religious, the latter of which Dr. Heegaard stubbornly treats
as synonymous with the religious in general, just as if there were only
one positive religion in one single confession. But where does the
religious come from that drops like a bomb into the book's conclusion? Is
it enough for a cognition of principles that we do not mistake subjective
for objective, religious for rational, without these concepts being seen in
the least as determinations of principle, as utterances of one and the same
being? The author says merely that revealed religion is manifestly a
% The Danish prints no stop after `Aandsmagt'; the English supplies it.
spiritual power. But what does that prove? For the cognition of principles
hardly much.

And what does it prove that the author calls the religious an ambiguous
% `i rationel og antirationel Forstand': Forstand here is `sense', not the
% faculty (understanding).
determination, since it can be taken both in a rational and in an
anti-rational sense; for where does the anti-rational come from? Why, it
comes from revelation. And revelation? Why, it is simply there. But is it
therefore anti-rational? Surely there is so much thought and reason in the
Old and New Testaments that it at least holds its own against the
miraculous. If one wants to draw the religious into the development as a
definite positive form of religion, one must accept the consequences and at
the same time draw it into the investigation. If one will not do this, but
proceeds from it in an entirely immediate way, then let it be done
properly, and let us have a catechism! That is incomparably easier than
these intricate expositions, which after all can only assert, not
demonstrate. But here the dualism has broken out that is now pressing hard
into our philosophical condition, and that shows this condition to be
neither more nor less than a crisis. It weakens the religious, which, as
% --- p. 31 ---
\opage{31}unconditionally contrary to reason, becomes an impotent spiritual
force, and it weakens the commands of reason, which are robbed of their
deepest elements in order to transfer these to the anti-rational, which
means little else than to shelve them. A so-called rational ethics that
cannot and dares not treat in its own way the ethical questions that every
revelation treats in its own way is without generative power and of no
significance; at most it becomes an ethics that ``can be applied in civil
society and the state as they are in reality,'' but not by a person in whom
the ethical, in whatever form, stirs with vigor.

The older principle of division is brought by the author into connection
with the new one by designating subjective-objective ethics as the one that
is conscious that it is theory and not practice, whereas the
representatives of merely objective and merely subjective ethics are said
to be caught in the illusion that theory and practice simply coincide ---
a strong expression indeed, when one recalls that among those who thus err
are thinkers such as Kant and Fichte. But to this new opposition of theory
and practice, this third principle of division, which like the second drops
from heaven at the end of the book, not the slightest attention has been
paid in what went before. One must, then, think the earlier forms through
again from this point of view and so do the work the author ought to have
done. He says himself that from this point of view one has greater sympathy
for Aristotle, the Stoics and the Epicureans than for Hegel and Fichte, but
why merely sympathy? The whole arrangement becomes a different one; Hegel
and Aristotle no longer come into the same class at all. On the relation
between the two principles of division that appeared last, moreover, Dr.~Heegaard
does not express himself with any clarity at all; one sees, to be
sure, that practical ethics can be both rational and religious; but for the
former there appears no independent representative; it is found only as a
stage of development under R. Nielsen.

% The Danish prints `etisko' for `etiske'; rendered as the sense.
The author has thus really wanted to view the ethical forms under a
threefold point of view, namely partly according to the opposition
subjective-objective, partly according to that of rational-religious, and
partly according to that of theoretical-practical. With this, however,
three principles of division have been set up. Now however unclear the
articulation might be, these three could nevertheless very well have been
united by virtue of dialectic. Under each of them all the forms would have
had to be capable of being
% --- p. 32 ---
\opage{32}placed; they would have had to interpenetrate at every point. It
ought to have been shown how the subjective is the objective, that both are
the religious, both the existential, the existential the theoretical, etc.,
until the whole cycle had been set in motion. Then the author could have
made his distinctions with a clear conscience. As it is, the forms bristle
outward away from one another without inner coherence, without any
dialectic at all, and yet the author's standpoint is the dialectical one,
so dialectical that he even --- something at which Kierkegaard would marvel
--- lets the very boundary between the rational and the anti-rational be of
a dialectical nature.

One question forces itself upon us in conclusion, a very important one,
namely this: Does the author himself have any standpoint? Because he, e.g.,
presents Feuerbach as a rational-atheistic ethicist who has no ethics, it
does not of course follow that the author is one. And because R. Nielsen is
presented as the one who gives a religious and rational ethics, one still
cannot presuppose it to be the author's assumption that such an ethics has
validity.

What would have to show this is missing, namely a critique of Nielsen's
view. For the remarkable fact is that when the author comes to the points
that, as the decisive ones, demanded a critique in the highest degree, he
very modestly closes his critical knife, the one he has used so vigorously
and bravely in what went before. Professor Nielsen is beyond comparison the
most significant representative in the whole last class, and the author
has him play a strange role. He is the representative of a distinct form,
but as such he is representative of representatives of other forms that
have been given no voice of their own, e.g., of the rational-religious.
Further, he is a critique of the whole and so also of himself, and finally
he is the chief representative of ethics, although he has written no
ethics.

What would one say if an art critic were lecturing on a sculptural group
rich in figures and, after giving his judgment of the less important
figures, vanished when he came to judge the principal figure, which cast
the real light over the whole? Could he excuse himself by saying that the
principal figure had to speak for itself and that he would say nothing but
what the principal figure clearly said? No, then he might have spared
himself the whole lecture; for if his listeners are able to judge the
principal figure, they can also judge
% --- p. 33 ---
\opage{33}the whole group. Either one gives a lecture, and then one gives it
to the end, or one leaves it to others to work the matter out for
themselves, and then there is no lecture at all.

If, then, it is Heegaard's opinion that Nielsen's standpoint is his own,
he must say so; it must not be a matter about which one knows anything only
from a more or less reliable quarter. What view underlies the critique is
at present a secret. Only this is certain: if the author has no other
standpoint than Professor Nielsen's, then it is just as strange to see
R.~Nielsen's name among those of the other thinkers as it would be to see the
author's own; if, on the other hand, he is not at one with Professor
Nielsen, he has hardly shown great independence or self-reliance by
leaving the latter's view unjudged.
\piecerule

% --- p. 34 ---
\piecehead{Rejoinder to Dr. Heegaard}{\opage{34}\emph{Rejoinder to Dr. P. S. V. Heegaard}\\[2pt](1867)}
In an anticritique whose tone is by no means justified by the
tone of the review against which it is directed, Dr. P. S. V.
Heegaard has not disdained to make use of the old trick with
which a man who, in an intellectual contest, has good reasons for not
wanting a reply sometimes tries to gag his opponent, and whose
recipe runs as follows: ``Should Mr. N. N. now, as would be just
like him, feel an urge to have the last word, he is free to do
so, etc.'', which, being interpreted, means: ``Since for my part I
have no other defense than to put forward denials and to throw
out accusations that cannot stand a moment longer than it takes
a reply to arrive, I feel a heartfelt urge to cut my opponent
off, if at all possible.''

Now, though it pains me not to be able to satisfy so deeply felt
an urge, I nonetheless beg the Doctor to be assured that I fully
understand and appreciate it. It is quite understandable that
Dr. Heegaard, feted and acclaimed by such expert judges as those
he has found in \emph{Illustreret Tidende} and \emph{Berlingske}, should
himself have begun to regard his work as good and thorough, and
that it should have been an unpleasant surprise to him to meet a
critic who was of the opposite opinion; but he ought to have taken
this fact more calmly; he ought to have made an effort to grasp
the line of thought in my article, at least where it lies right
under his nose, and he would then not have stopped in the middle
of reading a compound sentence to lament the meaninglessness of
its first half.

% --- p. 35 ---
\opage{35}Here is an example of how the Doctor has read. My review
contained the following words: ``A question forces itself upon us:
Does the author himself have any standpoint? Because he, e.g.,
presents Feuerbach as a rational-atheistic ethicist who has no ethics,
it does not of course follow that the author is one. And because
R. Nielsen is presented as the one who gives a religious and
rational ethics, one still cannot presuppose it to be the
author's assumption that such an ethics has validity.''

Dr. Heegaard now tears the one sentence about Feuerbach out of its
context, cannot grasp the point of this statement, since no
reasonable person could think of reasoning in that way, and cries:
``To what purpose, then, the remark?'' Now, if my honored opponent
really was unable, in a calm state of mind, to see what purpose
this remark serves, it would be in vain for me to try to explain
the matter to him; and it is only because, while writing down his
piece, he was evidently in a state of muddled irascibility that I
spoon-feed him the necessary enlightenment: the two cases set out
(that of Feuerbach and that of Nielsen) are of the same kind; if
nothing can be inferred from the first, then nothing can be
inferred from the second either, and consequently it is left
unexplained in the book itself whether the author is \emph{with}
Professor Nielsen or \emph{against} Professor Nielsen, or whether he
has no opinion at all.

In his book the Doctor showed himself an accomplished diplomat,
took careful care not to hint at any independent view, and
performed the egg dance of departing from Professor Nielsen on no
point whatever, without, however, hinting by so much as a syllable
at any agreement with him; in short, he conducted himself as a
cautious man.

Now caution is no doubt a priceless virtue, a burgomaster's virtue,
as they say --- and who knows whether Dr. Heegaard has not mistaken
the calling of his virtues by insisting on being a philosopher! ---
but there do occur cases here in life in which one must sacrifice
it. It needs only someone to prove so importunate as to ask whether
the Doctor, who ``is in on'' the development of others, has then no
development of his own at all, and he is forced to confess that he
has so far never thought of anything but the old thing, swearing by
the master's words; indeed, he affects to be almost offended at the
suspicion that he himself might have a view, and derives his
receptiveness toward Professor Nielsen
% --- p. 36 ---
\opage{36}``simply'' from the fact that one does not criticize the
standpoint on which one stands oneself\fnmark{*)}{Dr. Heegaard speaks of ``being in on the development'' much as the
Cat and the Hen in Andersen speak of purring and laying eggs. He seems
to think that one must, as a \emph{studiosus perpetuus}, have heard all of
Professor Nielsen's lectures, those for the freshmen and those for the
educated as well, if one is to be allowed to join in the talk about
philosophical subjects here at home. Indeed, he even invites me --- not
without wit --- to check whether he has correctly reproduced Professor
Nielsen's line of thought from lectures that have as yet been neither
delivered nor printed.}. But where in the world was I supposed to learn
anything about Dr. Heegaard's standpoint, when the only road to it
ran through an inference of which the Doctor himself (rashly)
declares that no reasonable person could think of drawing it?

No, the Doctor's newspaper article supplies a piece of information
of which his book gives not the slightest hint, and in imparting it
Dr. Heegaard has exchanged his arrogance for a modesty without
equal. Far be it from me to accuse my opponent of ``straining for
effect''; but might not the modesty he has here displayed produce
its effect all the same? Might not many a father or guardian say to
his ward: ``O my son! Look at yourself in the mirror of this mature
man! Though no longer quite young, though a Doctor of Philosophy and
a Privatdocent, though the author of a book that inspires awe by its
thickness alone, he is still so modest as to have no other opinion
than the one his teacher has.'' Oh, surely the Doctor, without
having intended it, will achieve such an effect, and his modesty
will propagate itself from generation to generation.

Now modesty is no doubt an inestimable virtue, in some people's
opinion even the finest in a young man, but there are nonetheless
cases here in life in which an expression of modesty becomes one
and the same thing as a declaration of bankruptcy. The Doctor is in
a great and lamentable error when he thinks that one does need a
standpoint of one's own to write an ethics but by no means to write
an introduction to one, and in another error just as great when he
supposes that my complaint about his lack of a standpoint has its
ground in a confusion of the introduction to ethics with ethics
itself.

Dr. Heegaard is bold enough to reject my definition
% --- p. 37 ---
% The Danish prints a point for a comma in `indrømme. at' (below); the
% English renders the sense.
% `Enten faar han en almindelig Metafysik, der falder udenfor de særskilte
% Videnskabers': construed as "a general metaphysics, falling outside that
% of the particular sciences", i.e. the introduction would become general
% metaphysics, which belongs to no special science.
\opage{37}of what an introduction to ethics means with a simple: This is
incorrect. We shall see with what right. It is already peculiar that
my opponent should insist on one particular explanation of the title
of his book, just as if ``Introduction to Ethics'' were a name for a
science current and accepted among men of science. But it is still
more peculiar that so cautious an author is not more cautious with
the definition of the concept that he gives. When he says that one
cannot develop what ethics is without explaining what the Good is,
this is self-evident, understood absolutely; but taken relatively it
is wrong, for if no distinction is made between the science of a
science and that science itself, the system of the sciences turns
into an even and uniform fog. And when he holds that the
introduction to ethics ought to consist in the metaphysics of
ethics, he runs aground on the following dilemma: either he gets a
general metaphysics, one that falls outside that of the particular
sciences, which is nonsense, or he is compelled to admit that the
metaphysics of the special discipline is not an independent science
but must consist of propositions borrowed from general metaphysics.
And that is really how matters stand; that is why, e.g., the
metaphysics of the beautiful has never been set up as a separate
science either. The most amusing thing, however, is that Dr.
Heegaard, by virtue of his way of seeing things, must consistently
demand metaphysics itself as philosophical propaedeutics, whereby he
comes into conflict partly with common sense, which is of lesser
importance here, and partly with his own oracle, Professor Nielsen,
who, in agreement with my conception, defines propaedeutics as the
science that answers the question: ``What is philosophy?''

But I will make Dr. Heegaard a present of his bold assertions on
this point; let him define the introduction as he pleases, this much
still stands unshaken: whatever an introduction may be, it must be an
organic component of the system to which it leads and by which it is
comprehended, and then the same life, the same spirit, the same
principle must also pulse in the introduction as in the system, and a
lack of principle must be just as reprehensible in the former as in
the latter. How, then, can so modest a man as Dr. Heegaard believe
that, with his modest lack of a standpoint, he is fit to write the
introduction?

Or does Dr. Heegaard perhaps have a principle after all, a
% --- p. 38 ---
\opage{38}principle without a standpoint? What, then, has he set against
my specific demonstration of his spineless wavering between three
different principles of division? Not a word, not a letter!
% The Danish prints `Saa frejdig kan ellers er' with `kan' for `han';
% the English renders the sense.
Bold as he otherwise is at denying, here he fails to hear the criticism, and
is wise to do so. And is his talk about the subdivision perhaps more
successful than his silence about the main point? Can anything more
preposterous be imagined than that the Doctor should call on \emph{me}
to show what compelling reason there is to subdivide the ethics of
essentiality, i.e., to break the phenomenon on the very wheel on
which he himself has laid it?

I answer: If there lies in the concept of actuality a compelling
reason for subdivision into ``immediate'' and ``reflected'' --- and
the Doctor must surely assume so, since he divides in that way ---
then such a reason must also lie in essentiality. For in the opposite
case those two determinations would, of course, fall outside the
latter --- a fine ontology. The emphasis lies on what is a matter of
principle, on the whole of all the determinations of principle
employed, which must form a unity. If the historical material does
not correspond, so much the worse for the Doctor. It is he, not I,
who has to make the material and the principles essentially coincide;
it is he who has registered, having been unable to show the
transition between the historical and what is a matter of principle,
and he will hardly succeed in making up for his short weight with the
ingenuous declaration that he ``has not, after all, given so little
in the way of development of principles.'' It does not become any
more certain because he himself assures us of it.

But then Dr. Heegaard has, of course, declared that a criticism of
his work which, like mine, does not take up all the questions afresh
from the beginning and does not carry out the insurmountable task of
putting the correct in the place of the supposedly incorrect --- such
a criticism ``is not worth heeding,'' and this pronouncement may have
made an impression on one or two. But how can an author who by his
own admission has a view of the whole only at second hand demand of a
newspaper review of two and a half columns so much more than he
himself has given in four hundred and fifty pages? The reviewer has
supplied nothing but assertions, he says, without proof. That is only
half true, but let it be so! Against these he sets other assertions,
at least as unproved,
% --- p. 39 ---
\opage{39}apparently, to be sure, without being aware that he offers
nothing more.

But then the question still remains which assertions the thinking
reader of the book is compelled to grant, his or mine, and I
confidently leave it to those who know the subject to judge between
us.
\piecerule

% --- p. 40 ---
\piecehead{The Logic of the Fundamental Ideas}{\opage{40}\emph{R. Nielsen: Grundideernes Logik}\\[2pt]{\normalsize[The Logic of the Fundamental Ideas]}}
% `2den Del' (Part 2) is kept in Danish as part of the title Brandes cites.
To give an idea of the line of thought in \emph{Professor R. Nielsen's}
``\emph{Grundideernes Logik, 2den Del}'' in a couple of words cannot be
done. To deliver a criticism would be more difficult still. But this
much shall simply be said: it would be a disgrace to our fatherland if
a work such as this did not find numerous buyers among the large
public that consists of the older or younger pupils of the excellent
university teacher and writer to whom it is due. The volume most
recently published is written with a logical virtuosity, a sharpness
and definiteness of expression, an omnipresent liveliness and flight
in the presentation, that must awaken the strongest admiration and
that make the reading as attractive as it is instructive. Even the
reader who is comparatively unfamiliar with investigations of this
kind will find large sections that treat, in a manner generally
intelligible and at the same time new, original and full of genius,
problems to which ordinary reasoning constantly returns and which it
is never able to solve by its own powers. Such are, e.g., the
questions of the origin and meaning of languages, of the necessity of
the Word as the expression of the divine thoughts, of God's
foreknowledge and omniscience, of the creation of the world, of
pre-existence, etc. The work is shot through with a struggle,
testifying to an unmistakable superiority, against Madvig, against
Sibbern, against Martensen and others as representatives of one-sided
or antiquated points of view, and with attacks, launched from new
standpoints, on thinkers of the past such as Kant and Hume, attacks
that cast a peculiar light on the merits and the errors of these men.

% --- p. 41 ---
% `lyst i Kuld og Køn': the old formula for declaring a child legitimate,
% i.e. publicly acknowledged as one's own; rendered "solemnly declared
% legitimate".
\opage{41}Since this place is so little suited to a critical pronouncement,
we shall permit ourselves a remark on only one quite single point.
Since \emph{Grundideernes Logik} has not, in Professor Nielsen's
treatment, become a strictly delimited science, and since the author
shows a pervasive endeavor to take in as much as possible of the
given religious content of representation, intimations crop up here
and there of the author's well-known views on the boundaries of the
domain of reason against the religious, although the work as yet
bears only faint traces of the monstrous doctrines that Dr. Heegaard
in his latest writing has presented as belonging to Professor
Nielsen. The cornerstone of this whole doctrine, however, is the
Kierkegaardian conception of ``the Paradox,'' and if only one could
succeed in showing to what desperate, indeed insane, consequences
adherence to this standpoint must necessarily lead, then there is hope
that the whole building erected upon it will collapse without further
criticism. But a point of attachment for such an attempt the author
of the Logic has himself provided. In seeking to show that the Kantian
category \emph{the thing in itself} is a mere phantom of mist, he
conducts the proof as follows: ``Kant cannot avoid passing a judgment
on the \emph{thing in itself}. He must say of it that it \emph{is}. But
once it has been proved that the category of the understanding,
\emph{being}, must be determining for the \emph{thing in itself}, then
the same holds of all the other categories.''

Nothing can be clearer or more correct than this line of thought. But
how, then, is it possible that the brilliant author of \emph{Grundideernes
Logik} can fail to see that this very same line of proof can be
applied against the concept ``the Paradoxical,'' which he has
nevertheless solemnly declared legitimate. The philosopher of the
absurd may writhe as much as he likes; he must still confess that, if
he can say nothing else about the Paradox, this impenetrable,
incomprehensible thing, he must at least establish of it that it
\emph{is}. But being is a determination of reason. And if ``the
Paradox'' now belongs to a circle within which the concept ``being''
is valid, then it also belongs to a circle within which all the
determinations of reason are valid. But what does this imply? Nothing
less than the flagrant contradiction that the Paradoxical, as subject,
receives nothing but predicates determined by reason. If one wishes to
avoid it, there is no other way out than to fall silent, i.e., not
merely to shut oneself up in an isolation in which, although one
cannot immediately
% --- p. 42 ---
\opage{42}communicate oneself to others, one nonetheless understands
oneself --- but to fall silent unconditionally, to become dumb as the
fish is dumb, to cease to be a human being. This is the dilemma. If
one holds fast to ``the Paradoxical,'' one will gradually find oneself
carpentering together a system like the Ptolemaic system of astronomy,
in which the first false conjecture could only ever be upheld by
others just as desperate, in the same way as the paradox that
conflicts with reason has already made a ``rational paradox'' (!)
necessary. But then, too, one fine day the clear Copernican system of
reason will stand victorious without any struggle, solely by the
all-conquering power that lies in the holy simplicity of the true.
\piecerule

\end{document}
